% ============================================================================
\clearpage
\section*{v4.9.0-r61 Source-Only: A Frame Audit of the Open Problems OP1 to OP5}
\addcontentsline{toc}{section}{v4.9.0-r61 Source-Only: A Frame Audit of the Open Problems OP1 to OP5}
\label{sec:v490-r61-op-frame-audit}
% ============================================================================

\paragraph{Purpose.}
The author suspected that the open problems, like the main statements before them, were studied in mixed frame configurations.  This layer audits OP1 to OP5 against
Convention~F (v4.9.0-r22, r24), the class split of v4.9.0-r18/r19 (scale quantities by way~C, \(e^{k\eta}\); readout quantities by way~A, \(\gL\)) and rule~7 of the round
format.  It makes no prediction.  One computation (\texttt{llm/k101\_fenster\_eta.py}) is an observation, not a test.  Earlier layers are not rewritten.

\noindent\textbf{Source-only status.}
This v4.9.0-r61 layer is source-only.  It does not build a PDF, does not assign a new DOI,
and does not authorize a public release.  The v4.9.0 Version DOI
\href{https://doi.org/10.5281/zenodo.23534954}{10.5281/zenodo.23534954} and the Concept DOI
\href{https://doi.org/10.5281/zenodo.19210728}{10.5281/zenodo.19210728} are unchanged by it.

\subsection*{v4.9.0-r61 results}
\label{subsec:v490-r61-results}

\paragraph{Theoremlet v4.9.0-r61.1 (OP1).}
\begin{enumerate}[leftmargin=2em]
\item[(a)] The Tribonacci condition \(\gL=1/\bsym\) of Theorem~\ref{thm:tribonacci} equates a readout quantity of the matter frame with \(1/\bsym=\beta_T\), the offset
of the tachyon in its own measure (v4.9.0-r24.2); equivalently \(\cosh2\eta=e^\eta\), way~A equal to way~C.  It is not a frame-consistent fixed point.
\item[(b)] The Hopf value of Theorem~\ref{thm:bsym-hopf} reads \(\bsym=\sin(\theta_H/2)\) for a pure state measured from the photon pole, i.e.\ in the reference frame R
(v4.9.0-r24).  At \(\bsym^2=\tfrac13\) that pure state has populations \((\tfrac23,\tfrac13)\), while the matter frame of Convention~F has \((\tfrac34,\tfrac14)\) and \(|b|=\tfrac12\)
(v4.9.0-r17).  The same \(\bsym\) labels two different states.
\item[(c)] The measured relation \(\gL=\gBI/\gstring\) of Theorem~\ref{thm:symmetric-frame} is the Lorentz factor between the real projections of M and T and the ratio of
Bloch lengths \(|b_R|/|b_M|\); it is frame-consistent, but an input.
\item[(d)] The Hopf gap and the Tribonacci gap, the subject of OP4, compare the measured matter value with values from another frame or another class.
\item[(e)] Observation: under the reading of v4.9.0-r55, r57, r60 the lepton data fix \(\eta\in[0.611,0.647]\), i.e.\ \(\bsym\in[0.545,0.570]\) and \(\gL\in[1.844,1.961]\),
independently of \(\gBI/\gstring\).  The measured value lies inside; Hopf (\(2\)) and Tribonacci (\(1.839\)) lie just outside.  Across the window
\(\sin^2\theta_{23}\in[0.491,0.498]\).
\end{enumerate}

\paragraph{Theoremlet v4.9.0-r61.2 (OP2 to OP5).}
\begin{enumerate}[leftmargin=2em]
\item[(a)] OP2 is closed except (T4), which is OP1 (Theorem~\ref{thm:op2-T4-equals-op1}); (T4) reads \(\gL\) as a ratio of spectral unfolding scales and carries the same
class mixing as r61.1(a).
\item[(b)] OP4 as posed rests on r61.1(d).
\item[(c)] The Theta sum of OP5 (Definition~\ref{def:theta-sum-B}) converts zeros (frequencies) and cutoff (a length), scale quantities, with \(\gL\) instead of \(e^{k\eta}\).
The difference is \(0.5\%\) (\(e^\eta=1.854\), \(\gL=1.864\)) and does not explain the failure of the conjecture by \(1.003\).
\item[(d)] OP3 was not re-examined.
\end{enumerate}

\paragraph{Proof sketch.}
(a)--(d) of r61.1 and r61.2 by comparison with the rules of v4.9.0-r17 to r24; the populations by direct computation; (e) by script \texttt{llm/k101\_fenster\_eta.py}
(output \texttt{llm/k101\_out.txt}) on the models of v4.9.0-r57 with a grid of \(0.001\).  \(\square\)

\paragraph{Reading.}
\begin{enumerate}[leftmargin=2em]
\item The suspicion holds for OP1: two of its three values come from configurations that mix frames or classes; only the measured value is frame-consistent.  OP1 stays
open, but its target is the matter value itself, not a correction of the Hopf or Tribonacci value.
\item OP4 is ill-posed in its present form; only OP1 remains.  OP2 (T4) carries the same mixing.  In OP5 the mixing is present but not the cause of the failure.
\item r61.1(e) is not blind (the reading \(\varepsilon=\eta\) was chosen after the windows were known).  It shows that lepton mixing, read this way, selects the measured
frame offset over the Hopf and Tribonacci values.
\end{enumerate}

\subsection*{v4.9.0-r61 status in the three frames}
\label{subsec:v490-r61-frames}

\begin{center}
\small
\begin{tabular}{@{}>{\raggedright\arraybackslash}p{0.30\textwidth}>{\raggedright\arraybackslash}p{0.62\textwidth}@{}}
\toprule
result & M / T / P\\
\midrule
r61.1(a) & mixes M (\(\gL\)) and T (\(\beta_T=1/\bsym\))\\
r61.1(b) & mixes R (pure Hopf state) and M (mixed matter state); P is the degenerate point \(b=0\)\\
r61.1(c), (e) & M; the M--T factor is frame-consistent\\
r61.2 & as stated per item\\
\bottomrule
\end{tabular}
\end{center}

\subsection*{v4.9.0-r61 register entries}
\label{subsec:v490-r61-register}

\begin{center}
\small
\begin{tabular}{@{}>{\raggedright\arraybackslash}p{0.34\textwidth}>{\raggedright\arraybackslash}p{0.58\textwidth}@{}}
\toprule
statement & status after v4.9.0-r61\\
\midrule
Tribonacci fixed point (Theorem~\ref{thm:tribonacci}) & frame and class mixing; candidate for a status note\\
Hopf value of \(\bsym\) (Theorem~\ref{thm:bsym-hopf}) & reads an R-frame pure state as the M offset; candidate for a status note\\
OP4 (continuum-limit correction) & ill-posed in its present form\\
OP1 & open; target is the matter value itself\\
\(\bsym\) from lepton mixing (reading r55--r60) & \([0.545,0.570]\); measured inside, Hopf and Tribonacci outside (not blind)\\
\bottomrule
\end{tabular}
\end{center}

\subsection*{v4.9.0-r61 ledger and guard update}
\label{subsec:v490-r61-ledger-guard}

\[
R_{\rm op}^{v490\text{-}r61}=(1,1,0)
\]
for (frame mixing in OP1, OP2 (T4), OP4 located; OP5 mixing found but not causal; OP1 itself open).

The layer spends no guard:
\[
R_{\rm glob}^{(v490\text{-}r61,+)}=(0,1,1,1,1),
\qquad
\Gamma_{\rm guard}^{v490\text{-}r61}=(1,1,1,1,1),
\qquad
\beta_4=1/256\ \text{unspent}.
\]
It does not imply exact finite \(C5\) zero closure, exact all-readout terminality,
\(C6/T2\), a global photon-clock/RH theorem, or global all-readout closure.
